%%%%%%%%%%%%%%%%%%%%%%%%%%%%%%%%%%%%%%%%%%%%%%%%%%%%%%%%%%%%%%%%%%%%%%
% Lớp 11 | Giới hạn. Hàm số liên tục | Giới hạn của hàm số
% Dạng 1D3?2-4: PP đặt thừa số chung, kết quả vô cực
% Số câu: 80  (sinh tự động bởi tach_cau.py; từ nay thầy sửa trực tiếp file này)
%%%%%%%%%%%%%%%%%%%%%%%%%%%%%%%%%%%%%%%%%%%%%%%%%%%%%%%%%%%%%%%%%%%%%%

%%%%%%%%%%%%%%%% TRẮC NGHIỆM 4 PHƯƠNG ÁN %%%%%%%%%%%%%%%%

% nguồn: Ok_Toan11_TracNghiem_Lan2
\begin{ex}%[1D3N2-4]%[ma:74501d04]
	Tìm $\lim\limits_{x\to-\infty} x^{2024}$ ta được  kết quả là
	\choice 
		{\True  $+\infty$}
		 {$0$}
		{$-\infty$} 
		 {$2024$}
	 \loigiai{
		Vì $2024$ là số nguyên  dương chẵn nên $\lim\limits_{x\to-\infty} x^{2024}=+\infty$.
	} 
\end{ex}

% nguồn: Ok_Toan11_TracNghiem_Lan2
\begin{ex}%[1D3N2-4]%[ma:0fe65317]
	 Tìm $\lim\limits_{x\to-\infty} x^{2023}$ ta được kết quả là
	 \choice
		 {$2023$}
		 {$0$}
		 {$+\infty$}
		{\True $-\infty$} 
	\loigiai{ 
		Vì $2023$ là số  tự nhiên lẻ nên $\lim\limits_{x\to-\infty} x^{2023}=-\infty$.
	} 
\end{ex}

% nguồn: Ok_TracNghiemLop11_Lan1
\begin{ex}%[1D3N2-4]%[ma:64bab882]
	Tính $\lim\limits_{x\to +\infty  } (-x^3+x^2+2)$.
	 \choice
		 {$0$}
		{$-\infty  $ }
		{\True $+\infty $ } 
		 {$2$}
	 \loigiai{
		$\lim\limits_{x\to +\infty }( -x^3+x^2+2)=\lim\limits_{x\to +\infty }\left[\left( -x^3 \right)\left( -1+\dfrac{1}{x}+\dfrac{2}{x^3} \right) \right]=\lim\limits_{x\to +\infty }\left( -x^3 \right)\cdot\lim\limits_{x\to  +\infty }\left( -1+\dfrac{1}{x}+\dfrac{2}{x^3} \right)$. \\
		Ta có: $\lim\limits_{x\to +\infty }\left( -x^3 \right)=-\infty $  và $\lim\limits_{x\to +\infty }\left( -1+\dfrac{1}{x}+\dfrac{2}{x^3} \right)=-1$.
		\\ 
	Vậy $\lim\limits_{x\to +\infty}(-x^3+x^2+2)=+\infty$.} 
\end{ex}

% nguồn: Ok_Toan11_TracNghiem_Lan2
\begin{ex}%[1D3H2-4]%[ma:8112dd27]
	Tìm  $\lim\limits_{x\to-\infty}(x^3+3x^2)$ ta được kết quả là
	\choice 
		{$0$} 
		{\True $-\infty$} 
		{$+\infty$} 
		{$1$} 
	\loigiai{ 
		 Ta có $\lim\limits_{x\to-\infty}(x^3+3x^2)=\lim\limits_{x\to-\infty} x^3\left(1+\dfrac{3}{x}\right)$.\\
		Ta có  $\lim\limits_{x\to-\infty} x^3=-\infty$ và $\lim\limits_{x\to-\infty}\left(1+\dfrac{3}{x}\right)=1 > 0$ nên $\lim\limits_{x\to-\infty}(x^3+3x^2)=-\infty$.
	 }
\end{ex}

% nguồn: Ok_Toan11_TracNghiem_Lan2
\begin{ex}%[1D3H2-4]%[ma:2069d81d]
	Tìm  giới hạn $\lim\limits_{x\to 2}\dfrac{x+1}{(2-x)^4}$.
	 \choice
		 {\True $+\infty$}
		{$-\infty$} 
		 {$-2$}
		 {$1$}
	\loigiai{ 
		Ta có  $\lim\limits_{x\to 2}\dfrac{x+1}{(2-x)^4}=\lim\limits_{x\to 2}(x+1)\cdot \dfrac{1}{(2-x)^4}=+\infty$ (vì $\lim\limits_{x\to 2}(x+1)=3>0$ và $\lim\limits_{x\to 2}\dfrac{1}{(2-x)^4}=+\infty$).
	} 
\end{ex}

% nguồn: Ok_Toan11_TracNghiem_Lan2
\begin{ex}%[1D3H2-4]%[ma:be74ffc3]
	Tính $\lim\limits_{x\to  -\infty}(\sqrt{2x^2+x}+x)$?
	 \choice
		{\True $+\infty$} 
		{$-1$} 
		 {$-\infty$}
		{$0$} 
	\loigiai{ 
		Ta có $\lim\limits_{x\to -\infty}(\sqrt{2x^2+x}+x) =\lim\limits_{x\to -\infty}\left[\sqrt{x^2\left(2+\dfrac{1}{x}\right)}+x\right]$\\ \hspace*{1cm} 
		$=\lim\limits_{x\to -\infty}\left(-x\sqrt{2+\dfrac{1}{x}}+x\right) =\lim\limits_{x\to  -\infty}\left[x\left(-\sqrt{2+\dfrac{1}{x}}+1\right)\right]$.\\
	Vì $\lim\limits_{x\to -\infty}x=-\infty$  và $=\lim\limits_{x\to -\infty}\left(-\sqrt{2+\dfrac{1}{x}}+1\right)=1-\sqrt{2}<0$ \\ Vậy $\lim\limits_{x\to -\infty}(\sqrt{2\mathrm{x^2}+x}+x)=+\infty$.}
\end{ex}

% nguồn: Ok_Toan11_TracNghiem_Lan2
\begin{ex}%[1D3H2-4]%[ma:c319f0da]
	Tìm $\lim\limits_{x\to 0}\left(\dfrac{1}{x}-\dfrac{1}{x^2}\right)$. 
	\choice 
		 {1}
		{0} 
		 {\True $-\infty$}
		 {$+\infty$}
	\loigiai{ 
	$\lim\limits_{x\to 0}\left(\dfrac{1}{x}-\dfrac{1}{x^2}\right)=\lim\limits_{x\to 0}\dfrac{x-1}{x^2}=-\infty$ vì $\lim\limits_{x\to  0}(x-1)=-1<0$, $\lim\limits_{x\to 0}x^2=0$ và $x^2>0$ với mọi $x\ne 0$.}
\end{ex}

% nguồn: Ok_Toan11_TracNghiem_Lan2
\begin{ex}%[1D3H2-4]%[ma:da2f27ac]
	 Giá trị của $\lim\limits_{x\to+\infty}\dfrac{1-x+x^2}{x+4}$ là
	\choice 
		{$1$} 
		{$0$} 
		{\True $+\infty$} 
		{$-\infty$} 
	\loigiai{ 
		Ta có $\lim\limits_{x\to+\infty}\dfrac{1-x+x^2}{x+4}=\lim\limits_{x\to+\infty}\dfrac{x^2\left(\dfrac{1}{x^2}-\dfrac{1}{x}+1\right)}{x\left(1+\dfrac{4}{x}\right)}=\lim\limits_{x\to+\infty} x\cdot \dfrac{\dfrac{1}{x^2}-\dfrac{1}{x}+1}{1+\dfrac{4}{x}}=+\infty$\\ 
		 (Vì $\lim\limits_{x\to+\infty}x=+\infty$ và $\lim\limits_{x\to+\infty}\dfrac{1-\dfrac{1}{x}+\dfrac{1}{x^2}}{\dfrac{1}{x}+\dfrac{4}{x^2}}=1>0$).
	 }
\end{ex}

% nguồn: Ok_Toan11_TracNghiem_Lan2
\begin{ex}%[1D3H2-4]%[ma:8bc0667d]
	Tính giới hạn $\lim\limits_{x\to -\infty}\dfrac{1+x-x^2}{x}$. 
	 \choice
		 {$0$}
		 {\True $+\infty $}
		{$1$} 
		 {$-\infty $}
	 \loigiai{
	 $\lim\limits_{x\to -\infty}\dfrac{1+x-x^2}{x}=\lim\limits_{x\to -\infty}\dfrac{x^2\left(\dfrac{1}{x^2}+\dfrac{1}{x}-1\right)}{x}=\lim\limits_{x\to -\infty}\left[x\left(\dfrac{1}{x^2}+\dfrac{1}{x}-1\right) \right]=+\infty $.}
\end{ex}

% nguồn: Ok_Toan11_TracNghiem_Lan2
\begin{ex}%[1D3H2-4]%[ma:335e800d]
	Giới  hạn $\lim\limits_{x\to-\infty}\dfrac{2x^3+5x^2-3}{x^2+6x+3}$ có kết quả là
	 \choice
		{$-2$} 
		{$+\infty$} 
		{\True $-\infty$} 
		 {$2$}
	\loigiai{ 
		Ta có $\lim\limits_{x\to-\infty}\dfrac{2x^3+5x^2-3}{x^2+6x+3}=\lim\limits_{x\to-\infty}x\cdot \dfrac{2+\dfrac{5}{x}-\dfrac{3}{x^3}}{1+\dfrac{6}{x}+\dfrac{3}{x^2}}=-\infty$\\ (vì $\lim\limits_{x\to-\infty}x=-\infty$ và  $\lim\limits_{x\to-\infty}\dfrac{2+\dfrac{5}{x}-\dfrac{3}{x^3}}{1+\dfrac{6}{x}+\dfrac{3}{x^2}}=2>0$).
	 }
\end{ex}

% nguồn: Ok_Toan11_TracNghiem_Lan2
\begin{ex}%[1D3H2-4]%[ma:80004923]
	Biết $\lim\limits_{x\to +\infty}\dfrac{(2-a)x-3}{x-\sqrt{x^2+1}}=+\infty$ (với $a$ là  tham số). Giá trị nhỏ nhất của $P=a^2-2a+4$ là
	 \choice
		{\True  $4$}
		{$3$} 
		{$5$} 
		{$1$} 
	\loigiai{ 
		Ta có  $\lim\limits_{x\to +\infty}\dfrac{(2-a)x-3}{x-\sqrt{x^2+1}}=\lim\limits_{x\to +\infty}\left[-\left((2-a)x-3\right)(x+\sqrt{x^2+1})\right]$\\ \hspace*{1cm}
		$=\lim\limits_{x\to +\infty}x^2\left[-\left((2-a)-\dfrac{3}{x}\right)\left(1+\sqrt{1+\dfrac{1}{x^2}}\right)\right] =+\infty  \Rightarrow -(2-a)\ge 0 \Leftrightarrow a\ge 2$.\\
		Với $a\ge 2$ $\Rightarrow a(a-2)\ge 0$ suy ra  $P=a(a-2)+4\ge 4$.\\ Dấu \lq\lq =\rq\rq xảy ra khi $a$ = 2. Vậy $GTNN$ của $P$ là $4$.
	 }
\end{ex}

% nguồn: Ok_Toan11_TracNghiem_Lan2
\begin{ex}%[1D3H2-4]%[ma:c5fce22c]
	Cho $\lim\limits_{x\to 1}\dfrac{x^3-1}{x^2-1}=\dfrac{a}{b}$ với $a$, $b$ là các số nguyên dương và $\dfrac{a}{b}$  là phân số tối giản. Tính tổng $S=a+b$.
	\choice 
		 {\True$5$}
		 {$10$}
		{$3$} 
		{$4$} 
	\loigiai{ 
		 Ta có $\lim\limits_{x\to 1}\dfrac{x^3-1}{x^2-1}=\lim\limits_{x\to 1}\dfrac{x^2+x+1}{x+1}=\dfrac{3}{2}$.\\
		Suy ra $\heva{&a=3\\ 
		&b=2.}$\\ 
		 Vậy$S=5$.
	 }
\end{ex}

% nguồn: Ok_Toan11_TracNghiem_Lan2
\begin{ex}%[1D3H2-4]%[ma:8dca55f9]
	Giá trị của $\lim\limits_{x\to -\infty}\dfrac{2x-1}{\sqrt{x^2+1}-1}$ bằng 
	 \choice
		{0} 
		{\True $-2$} 
		{$-\infty$} 
		 {2}
	\loigiai{ 
	Ta có $\lim\limits_{x\to -\infty}\dfrac{2x-1}{\sqrt{x^2+1}-1}  =\lim\limits_{x\to -\infty}\dfrac{2x-1}{-x\sqrt{1+\dfrac{1}{x^2}}-1} =\lim\limits_{x\to -\infty}\dfrac{2-\dfrac{1}{x}}{-\sqrt{1+\dfrac{1}{x^2}}-\dfrac{1}{x}} =-2$.}
\end{ex}

% nguồn: Ok_Toan11_TracNghiem_Lan2
\begin{ex}%[1D3H2-4]%[ma:e2aeb846]
	Tìm $\lim\limits_{x\to  -\infty}(x^3-3x^2+2)$.
	 \choice
		{$+\infty$} 
		{\True  $-\infty$}
		 {1}
		 {3}
	\loigiai{ 
		Ta có $\lim\limits_{x\to -\infty}(x^3-3x^2+2)=\lim\limits_{x\to  -\infty}x^3\left(1-\dfrac{3}{x}+\dfrac{2}{x^3}\right)=-\infty$\\
	 vì $\lim\limits_{x\to -\infty}x^3=-\infty$ và $\lim\limits_{x\to -\infty}\left(1-\dfrac{3}{x}+\dfrac{2}{x^3}\right)=1>0$.}
\end{ex}

% nguồn: Ok_Toan11_TracNghiem_Lan2
\begin{ex}%[1D3H2-4]%[ma:7e9d1bd7]
	 Kết quả của $\lim\limits_{x\to-\infty} (4x^5-3x^3+x+1)$ là
	 \choice
		 {\True $-\infty$}
		 {$0$}
		{$4$} 
		 {$+\infty$}
	 \loigiai{
		Vì $\lim\limits_{x\to-\infty}\left(4-\dfrac{3}{x^2}+\dfrac{1}{x^4}+\dfrac{1}{x^5}\right)=4$ và $\lim\limits_{x\to-\infty}  x^5=-\infty$ nên
		\begin{align*}
			\lim\limits_{x\to-\infty} (4x^5-3x^3+x+1)=\lim\limits_{x\to-\infty}  x^5\left(4-\dfrac{3}{x^2}+\dfrac{1}{x^4}+\dfrac{1}{x^5}\right)=-\infty.
		\end{align*}
	 }
\end{ex}

% nguồn: Ok_Toan11_TracNghiem_Lan2
\begin{ex}%[1D3H2-4]%[ma:237c1edb]
	 Giới hạn $\lim\limits_{x\to +\infty}\dfrac{cx^2+a}{x^2+b}$ bằng
	 \choice
		{$a$} 
		 {$b$}
		 {\True $c$}
		 {$\dfrac{a+b}{c}$}
	 \loigiai{
	Ta có $\lim\limits_{x\to +\infty}\dfrac{cx^2+a}{x^2+b}=\lim\limits_{x\to +\infty}\dfrac{c+\dfrac{a}{x^2}}{1+\dfrac{b}{x^2}}=\dfrac{c+0}{1+0}=c$.} 
\end{ex}

% nguồn: Ok_Toan11_TracNghiem_Lan2
\begin{ex}%[1D3H2-4]%[ma:775975a9]
	 Tính giới hạn $\lim\limits_{x\to-\infty}\left(\sqrt{x^2-x+1}-x\right)$.
	 \choice
		{$0$} 
		 {\True $+\infty$}
		 {$-\dfrac{1}{2}$}
		{$-\infty$} 
	 \loigiai{
	 \[\lim\limits_{x\to-\infty}\left(\sqrt{x^2-x+1}-x\right)=\lim\limits_{x\to-\infty}\left[-x\left(\sqrt{1-\dfrac{1}{x}+\dfrac{1}{x^2}}+1\right)\right]=+\infty.\]
	 }
\end{ex}

% nguồn: Ok_Toan11_TracNghiem_Lan2
\begin{ex}%[1D3H2-4]%[ma:f055825f]
	Giới  hạn $\lim\limits_{x\to-2}\dfrac{x+1}{(x+2)^2}$ bằng
	 \choice
		{\True $-\infty$} 
		 {$\dfrac{3}{16}$}
		 {$0$}
		 {$+\infty$}
	 \loigiai{
		Ta có $\lim\limits_{x\to-2}\dfrac{x+1}{(x+2)^2}=\lim\limits_{x\to-2}\dfrac{1}{(x+2)^2}\cdot (x+1)=-\infty$ (vì $\lim\limits_{x\to-2}\dfrac{1}{(x+2)^2}=+\infty$ và $\lim\limits_{x\to-2}  (x+1)=-1 < 0$).
	 }
\end{ex}

% nguồn: Ok_Toan11_TracNghiem_Lan2
\begin{ex}%[1D3H2-4]%[ma:23dbf030]
	Biết $\lim\limits_{x\to +\infty}\dfrac{\left( 2-a \right)x-3}{x-\sqrt{x^2+1}}=+\infty $. Giá trị  nhỏ nhất của $P=a^2-2a+4$ là
	\choice 
		{\True  $4$}
		{$3$} 
		 {$5$}
		 {$1$}
	\loigiai{ 
		 Ta có
		\begin{eqnarray*}
			& &\lim\limits_{x\to +\infty}\dfrac{\left( 2-a \right)x-3}{x-\sqrt{x^2+1}}\\ 
			&=&\lim\limits_{x\to  +\infty}\left[-\left( \left( 2-a \right)x-3 \right)\left( x+\sqrt{x^2+1} \right) \right]=+\infty.
		\end{eqnarray*}
		 $\Rightarrow -\left( 2-a \right)\ge 0\Leftrightarrow a\ge 2$.\\
	Với $a\ge 2\Rightarrow a\left(  a-2 \right)\ge 0$ suy ra $P=a\left( a-2 \right)+4\ge 4$.}
\end{ex}

% nguồn: Ok_Toan11_TracNghiem_Lan2
\begin{ex}%[1D3H2-4]%[ma:f6afb348]
	Giới  hạn $\lim\limits_{x\to-\infty}\dfrac{2x^2-x+1}{x+2}$ bằng
	 \choice
		 {$+\infty$}
		 {\True $-\infty$}
		 {$-2$}
		{$1$} 
	 \loigiai{
		Ta  có $\lim\limits_{x\to-\infty}\dfrac{2x^2-x+1}{x+2}=\lim\limits_{x\to-\infty}x\cdot \dfrac{2-\dfrac{1}{x}+\dfrac{1}{x^2}}{1+\dfrac{2}{x}}=-\infty$.
	} 
\end{ex}

% nguồn: Ok_Toan11_TracNghiem_Lan2
\begin{ex}%[1D3H2-4]%[ma:a203edf1]
	Giá trị của giới  hạn $\lim\limits_{x\to +\infty}(\sqrt[3]{3x^3-1}+\sqrt{x^2+2})$ là
	\choice 
		{$\sqrt[3]{3}+1$} 
		{\True $+\infty$} 
		 {$\sqrt[3]{3}-1$}
		{$-\infty$} 
	\loigiai{ 
		 $\lim\limits_{x\to +\infty}(\sqrt[3]{3x^3-1}+\sqrt{x^2+2})=\lim\limits_{x\to +\infty}x\left(\sqrt[3]{3-\dfrac{1}{x^3}}+\sqrt{1+\dfrac{2}{x^2}}\right)=+\infty$\\
	vì $\lim\limits_{x\to +\infty}x=+\infty$ và $\lim\limits_{x\to  +\infty}\left(\sqrt[3]{3-\dfrac{1}{x^3}}+\sqrt{1+\dfrac{2}{x^2}}\right)=\sqrt[3]{3}+1>0$.}
\end{ex}

% nguồn: Ok_Toan11_TracNghiem_Lan2
\begin{ex}%[1D3H2-4]%[ma:2017c7eb]
	Chọn kết  quả đúng của $\lim\limits_{x\to -\infty}(-4x^5-3x^3+x+1)$.
	 \choice
		 {$0$}
		{\True  $+\infty$}
		{$-\infty$} 
		 {$-4$}
	\loigiai{ 
		Ta có $\lim\limits_{x\to -\infty}(-4x^5-3x^3+x+1)  =\lim\limits_{x\to -\infty}x^5\left(-4-\dfrac{3}{x^2}+\dfrac{1}{x^4}+\dfrac{1}{x^5}\right) =+\infty$.\\
	 Vì $\heva{&\lim\limits_{x\to -\infty}\left(-4-\dfrac{3}{x^2}+\dfrac{1}{x^4}+\dfrac{1}{x^5}\right)=-4<0 \\&\lim\limits_{x\to -\infty}x^5=-\infty.}$}
\end{ex}

% nguồn: Ok_Toan11_TracNghiem_Lan2
\begin{ex}%[1D3H2-4]%[ma:d1f1857a]
	Giới hạn $\lim\limits_{x\to -\infty}(3x-\sqrt{9x^2-1})$  bằng:
	 \choice
		{$+\infty$} 
		 {0}
		{\True  $-\infty$}
		 {$-1$}
	\loigiai{ 
		$\lim\limits_{x\to -\infty}(3x-\sqrt{9x^2-1})=\lim\limits_{x\to  -\infty}\left(3x+x\sqrt{9-\dfrac{1}{x^2}}\right)=\lim\limits_{x\to -\infty}x\left(3+\sqrt{9-\dfrac{1}{x^2}}\right)=-\infty$\\
	vì  $\lim\limits_{x\to -\infty}x=-\infty$ và $\lim\limits_{x\to -\infty}\left(3+\sqrt{9-\dfrac{1}{x^2}}\right)=6>0$.}
\end{ex}

% nguồn: Ok_Toan11_TracNghiem_Lan2
\begin{ex}%[1D3H2-4]%[ma:94ccddaf]
	Tìm $\lim\limits_{x\to-3}\dfrac{2x+5}{(x+3)^2}$ ta được kết quả là 
	\choice 
		{$0$} 
		{$+\infty$} 
		{\True  $-\infty$}
		{$2$} 
	 \loigiai{
		Ta  có $\lim\limits_{x\to-3} (2x+5)=-1 < 0$, $\lim\limits_{x\to-3} (x+3)^2=0$ và $(x+3)^2 > 0$ khi $x\to-3$.\\
		 Vậy $\lim\limits_{x\to-3}\dfrac{2x+5}{(x+3)^2}=-\infty$.
	 }
\end{ex}

% nguồn: Ok_Toan11_TracNghiem_Lan2
\begin{ex}%[1D3H2-4]%[ma:c63b852a]
	 Tính $\lim\limits_{x\to-\infty}\sqrt{2x^2+x+1}$.
	\choice 
		{\True $+\infty$} 
		 {$-1$}
		 {$-\infty$}
		{$0$} 
	\loigiai{ 
		Ta  có
		\allowdisplaybreaks 
		\begin{eqnarray*}
			\lim\limits_{x\to-\infty}\sqrt{2x^2+x+1}&=&\lim\limits_{x\to-\infty}\sqrt{x^2\left(2+\dfrac{1}{x}+\dfrac{1}{x^2}\right)}\\ 
			&=&\lim\limits_{x\to-\infty}|x|\cdot  \sqrt{2+\dfrac{1}{x}+\dfrac{1}{x^2}}\\
			 &=&\lim\limits_{x\to-\infty}(-x)\cdot \sqrt{2+\dfrac{1}{x}+\dfrac{1}{x^2}}=+\infty.
		\end{eqnarray*}
		(Vì $\lim\limits_{x\to-\infty}(-x)=+\infty$  và $\lim\limits_{x\to-\infty}\sqrt{2+\dfrac{1}{x}+\dfrac{1}{x^2}}=\sqrt{2}>0$).
	 }
\end{ex}

% nguồn: Ok_Toan11_TracNghiem_Lan2
\begin{ex}%[1D3H2-4]%[ma:2ae4a982]
	Tìm  $\lim\limits_{x\to -\infty}(\sqrt{x^2+x}+2x)$.
	 \choice
		 {$2$}
		{\True  $-\infty$}
		{$1$} 
		 {$+\infty$}
	 \loigiai{
		Ta có $\lim\limits_{x\to -\infty}(\sqrt{x^2+x}+2x)$  $=\lim\limits_{x\to -\infty}\left(|x|\sqrt{1+\dfrac{1}{x}}+2x\right)$ $=\lim\limits_{x\to -\infty}\left(-x\sqrt{1+\dfrac{1}{x}}+2x\right)$\\\hspace*{1cm}
		 $=\lim\limits_{x\to -\infty}\left[x\left(2-\sqrt{1+\dfrac{1}{x}}\right)\right]$ $=-\infty$.
	\\vì  $\lim\limits_{x\to -\infty}x=-\infty$ và $\lim\limits_{x\to -\infty}\left(2-\sqrt{1+\dfrac{1}{x}}\right)=1$.}
\end{ex}

% nguồn: Ok_Toan11_TracNghiem_Lan2
\begin{ex}%[1D3H2-4]%[ma:8cb09ef1]
	$\lim\limits_{x\to -\infty}\dfrac{2x^3+5x^2-3}{x^2+6x+3}$  có kết quả là
	\choice 
		{$-2$} 
		{$+\infty$} 
		 {\True $-\infty$}
		{$2$} 
	 \loigiai{
		 Ta có $\lim\limits_{x\to -\infty}\dfrac{2x^3+5x^2-3}{x^2+6x+3}
		=\lim\limits_{x\to -\infty}\left(x\cdot  \dfrac{2+\dfrac{5}{x}-\dfrac{3}{x^3}}{1+\dfrac{6}{x}+\dfrac{3}{x^2}}\right)
		=-\infty$. 
	 }
\end{ex}

% nguồn: Ok_Toan11_TracNghiem_Lan2
\begin{ex}%[1D3H2-4]%[ma:e966694b]
	Trong các mệnh đề sau, mệnh đề  nào đúng?
	 \choice
		{\True $\lim\limits_{x\to -\infty}\dfrac{\sqrt{x^4-x}}{1-2x}=+\infty $} 
		{$\lim\limits_{x\to  -\infty}\dfrac{\sqrt{x^4-x}}{1-2x}=1$}
		{$\lim\limits_{x\to -\infty}\dfrac{\sqrt{x^4-x}}{1-2x}=-\infty $} 
		 {$\lim\limits_{x\to -\infty}\dfrac{\sqrt{x^4-x}}{1-2x}=0$}
	 \loigiai{
		Ta  có $\lim\limits_{x\to -\infty}\dfrac{\sqrt{x^4-x}}{1-2x}=\lim\limits_{x\to -\infty}\dfrac{x^2\cdot \sqrt{1-\dfrac{1}{x^3}}}{x\left( \dfrac{1}{x}-2 \right)}=\lim\limits_{x\to -\infty}\dfrac{x\sqrt{1-\dfrac{1}{x^3}}}{\dfrac{1}{x}-2}=+\infty$.\\
		Vì $\lim\limits_{x\to  -\infty} x=-\infty$ và $\lim\limits_{x\to -\infty}\dfrac{\sqrt{1-\dfrac{1}{x^3}}}{\dfrac{1}{x}-2}=-\dfrac{1}{2}$.
	 }
\end{ex}

% nguồn: Ok_Toan11_TracNghiem_Lan2
\begin{ex}%[1D3H2-4]%[ma:07e2eb62]
	Kết quả của $\lim\limits_{x\to+\infty}\dfrac{2x^3+5x^2-3}{x^2+6x+3}$ là 
	 \choice
		{$-2$} 
		 {\True $+\infty$}
		 {$-\infty$}
		 {$2$}
	 \loigiai{
		Ta có $\lim\limits_{x\to+\infty}\dfrac{2x^3+5x^2-3}{x^2+6x+3}=\lim\limits_{x\to+\infty}  x\cdot \dfrac{2+\dfrac{5}{x}-\dfrac{3}{x^3}}{1+\dfrac{6}{x}+\dfrac{3}{x^2}}=+\infty$.
	} 
\end{ex}

% nguồn: Ok_Toan11_TracNghiem_Lan2
\begin{ex}%[1D3H2-4]%[ma:c8c6af39]
	Giới hạn $\lim\limits_{x\to-\infty}\left(\sqrt{x^2-4x}-x\right)$  bằng
	 \choice
		 {$-2$}
		 {$-\infty$}
		{\True $+\infty$} 
		 {$0$}
	\loigiai{ 
	$\lim\limits_{x\to-\infty}\left(\sqrt{x^2-4  x}-x\right)=\lim\limits_{x\to-\infty}(-x)\left(\sqrt{1-\dfrac{4}{x}}+1\right)=+\infty$.
	} 
\end{ex}

% nguồn: Ok_Toan11_TracNghiem_Lan2
\begin{ex}%[1D3H2-4]%[ma:ee4a2ef9]
	 Tính $\lim\limits_{x\to -\infty}\dfrac{2x-3}{\sqrt{x^2+1}-x}$?
	\choice 
		{$0$} 
		{$-\infty$} 
		{\True $-1$} 
		 {$1$}
	\loigiai{ 
		 Ta có $\lim\limits_{x\to -\infty}\dfrac{2x-3}{\sqrt{x^2+1}-x}=\lim\limits_{x\to -\infty}\dfrac{2x-3}{\sqrt{x^2(1+\dfrac{1}{x^2})}-x}=\lim\limits_{x\to -\infty}\dfrac{2x-3}{-x\sqrt{1+\dfrac{1}{x^2}}-x}$
	$=\lim\limits_{x\to  -\infty}\dfrac{2-\dfrac{3}{x}}{-\sqrt{1+\dfrac{1}{x^2}}-1}=-1$.}
\end{ex}

% nguồn: Ok_Toan11_TracNghiem_Lan2
\begin{ex}%[1D3H2-4]%[ma:e2a35049]
	Để $\lim\limits_{x\to -\infty}\dfrac{\sqrt{4x^2+x+1}+4}{mx-2}=\dfrac{1}{2}$. Giá trị của $m$  thuộc tập hợp nào sau đây?
	\choice 
		 {$\left[3;6\right]$}
		{$\left[-3;0\right]$} 
		 {\True $\left[-6;-3\right]$}
		 {$\left[1;3\right]$}
	 \loigiai{
		Khi $m=0$ $\lim\limits_{x\to -\infty}\dfrac{\sqrt{4x^2+x+1}+4}{-2}=\lim\limits_{x\to -\infty}x\cdot \dfrac{\sqrt{4+\dfrac{1}{x}+\dfrac{1}{x^2}}+\dfrac{4}{x}}{2}=-\infty$ nên $m=0$ không thỏa  mãn.\\
		Khi $m\ne 0$,  ta có $\lim\limits_{x\to -\infty}\dfrac{\sqrt{4x^2+x+1}+4}{mx-2} =\lim\limits_{x\to -\infty}\dfrac{-\sqrt{4+\dfrac{1}{x}+\dfrac{1}{x^2}}+\dfrac{4}{x}}{m-\dfrac{2}{x}} =-\dfrac{2}{m}$.\\
	Theo bài ra ta có $-\dfrac{2}{m}=\dfrac{1}{2} \Leftrightarrow  m=-4\in \left[-6;-3\right]$.}
\end{ex}

% nguồn: Ok_Toan11_TracNghiem_Lan2
\begin{ex}%[1D3H2-4]%[ma:eac0ba7c]
	Tìm $\lim\limits_{x\to  1}\dfrac{x^4-3x^2+2}{x^3+2x-3}$.
	\choice 
		 {$-\dfrac{5}{2}$}
		 {\True$-\dfrac{2}{5}$}
		 {$\dfrac{1}{5}$}
		{$+\infty$} 
	 \loigiai{
		Ta có 
		$\lim\limits_{x\to 1}\dfrac{x^4-3x^2+2}{x^3+2x-3}=\lim\limits_{x\to 1} \dfrac{(x-1)(x+1)\left(x^2-2\right)}{(x-1)\left(x^2+x+3\right)}=\lim\limits_{x\to 1}  \dfrac{(x+1)\left(x^2-2\right)}{x^2+x+3}=-\dfrac{2}{5}$.
	} 
\end{ex}

% nguồn: Ok_TracNghiemLop11_Lan1
\begin{ex}%[1D3H2-4]%[ma:18f72928]
	Giá trị của giới hạn $\lim\limits_{x\to+\infty}\left(\sqrt{x^2+1}+x\right)$ là 
	 \choice
		 {$0$}
		{\True $+\infty$} 
		{$\sqrt 2-1$} 
		 {$-\infty$}
	\loigiai{ 
		Ta  có $\lim\limits_{x\to+\infty}\left(\sqrt{x^2+1}+x\right)= \lim\limits_{x\to+\infty}x\left(\sqrt{1+\dfrac{1}{x^2}}+1\right)= +\infty$.\\
		 vì $\heva{
		&\lim\limits_{x\to+\infty}x=+\infty\\ 
		 &\lim\limits_{x\to{+\infty}}\sqrt{1+\dfrac{1}{x^2}}+1=2 > 0.}$
	} 
\end{ex}

% nguồn: Ok_TracNghiemLop11_Lan1
\begin{ex}%[1D3H2-4]%[ma:14187887]
	Giá trị  của giới hạn $\lim\limits_{x\to+\infty}\left(\sqrt{1+2x^2}-x\right)$ là
	\choice 
		 {$0$}
		{\True $+\infty  $}
		{$\sqrt 2-1$} 
		{$-\infty $} 
	\loigiai{ 
		Ta  có $\lim\limits_{x\to+\infty}\left(\sqrt{1+2x^2}-x\right) =\lim\limits_{x\to+\infty}x\left(\sqrt{\dfrac{1}{x^2}+2}-1\right) =+\infty$\\
		Vì $\lim\limits_{x\to+\infty}x=+\infty$,  $\lim\limits_{x\to+\infty}\left(\sqrt{\dfrac{1}{x^2}+2}-1\right)=\sqrt 2-1 > 0$.
	 }
\end{ex}

% nguồn: Ok_TracNghiemLop11_Lan1
\begin{ex}%[1D3H2-4]%[ma:e7099891]
	Giá trị của giới hạn $\lim\limits_{x\to+\infty}\left(\sqrt[3]{3x^3-1}+\sqrt{x^2+2}\right)$  là
	 \choice
		 {$\sqrt[3]{3}+1$}
		{\True  $+\infty$}
		{$\sqrt[3]{3}-1$} 
		 {$-\infty$}
	\loigiai{ 
		Ta  có
		 $\lim\limits_{x\to+\infty}\left(\sqrt[3]{3x^3-1}+\sqrt{x^2+2}\right)= \lim\limits_{x\to+\infty}x\left(\sqrt[3]{3-\dfrac{1}{x^3}}+\sqrt{1+\dfrac{2}{x^2}}\right)=+\infty$.
	 }
\end{ex}

% nguồn: Ok_TracNghiemLop11_Lan1
\begin{ex}%[1D3H2-4]%[ma:b62bcd56]
	Một cấp số nhân có hai số hạng liên  tiếp là $16$ và $36$. Số hạng tiếp theo là
	\choice 
		{$720$} 
		{\True  $81$}
		 {$64$}
		 {$56$}
	\loigiai{ 
		Ta có cấp số nhân $\left(u_n\right)$ có 
		$ 
		\heva{ 
		& u_k=16  \\
		 & u_{k+1}=36
		} \Rightarrow q=\dfrac{u_{k+1}}{u_k}=\dfrac{9}{4}  \Rightarrow u_{k+2}=u_{k+1} q=81
		$. 
	} 
\end{ex}

% nguồn: Ok_Toan11_TracNghiem_Lan2
\begin{ex}%[1D3V2-4]%[ma:c820f1e6]
	Tính $\lim\limits_{x\to+\infty} \left(3x-2-\sqrt{4x^2-4x-3} \right)$. 
	 \choice
		 {$-1$}
		 {\True $+\infty$}
		{$-\infty$} 
		 {$5$}
	 \loigiai{
		Ta có $\lim\limits_{x\to+\infty}  \left(3x-2-\sqrt{4x^2-4x-3} \right)=\lim\limits_{x\to+\infty} x\cdot\left(3-\dfrac{2}{x}-\sqrt{4-\dfrac{4}{x}-\dfrac{3}{x^2}} \right)=+\infty$.\\
		Vì $\lim\limits_{x\to+\infty}x=+\infty$ và  $\lim\limits_{x\to+\infty}\left(3-\dfrac{2}{x}-\sqrt{4-\dfrac{4}{x}-\dfrac{3}{x^2}} \right)=3>0$.
	 }
\end{ex}

% nguồn: Ok_TracNghiemLop11_Lan1
\begin{ex}%[1D3K2-4]%[ma:cf8e8a5f]
	Xét tính tăng giảm và bị chặn của dãy số $\left\{\begin{aligned}
	 &u_1=2\\
	 &u_2=3\\
	& u_{n+1}=\sqrt{u_n}+\sqrt{u_{n-1}}, \,\forall  n \geq 2.\\
\end{aligned}\right. $
\choice 
	 {\True Tăng, bị chặn}
	{Giảm, bị  chặn}
	{Tăng,  chặn dưới}
	{Giảm, chặn trên} 
 \loigiai{
	Trước hết ta chứng minh $1<u_n<4, \,\forall  n$. \\
	Điều này hiển nhiên đúng  với $n=1$. \\
	 Giả sử $1<u_n<4$, ta có $1<u_{n+1}=\sqrt{u_n}+\sqrt{u_{n-1}}<\sqrt{4}+\sqrt{4}=4$. \\
	Ta chứng minh $(u_n)$ là  dãy tăng. \\
	Ta có $u_1<u_2$, giả sử $u_{n-1}<u_n, \,\forall n  \leq k$. \\
	Khi đó $\left\{\begin{aligned}
	 &u_k<u_{k-1} \\
	&  u_{k-1}<u_{k-2}\\
	\end{aligned}\right.
	\Rightarrow  \sqrt{u_k}+\sqrt{u_{k-1}}<\sqrt{u_{k-1}}+\sqrt{u_{k-2}} \Rightarrow u_{k+1}<u_k$. \\
	Vậy dãy $(u_n)$ là dãy tăng và bị  chặn.
	} 
\end{ex}

% nguồn: Ok_TracNghiemLop11_Lan1
\begin{ex}%[1D3B2-4]%[ma:ebcc17d1]
	Xét  tính bị chặn của dãy số $(u_n)$ với $u_n=4-3n-n^2$.
	 \choice
		 {Bị chặn}
		 {Không bị chặn}
		{\True Bị chặn trên} 
		{Bị chặn  dưới}
	\loigiai{ 
	Ta có $\forall n \in \mathbb{N}^*, u_n=\dfrac{25}{4}-\left(n+\dfrac{3}{2}\right)^2<\dfrac{25}{4} \Rightarrow (u_n)$ bị chặn trên; dãy  $(u_n)$ không bị chặn dưới.}
\end{ex}

% nguồn: Ok_TracNghiemLop11_Lan1
\begin{ex}%[1D3B2-4]%[ma:eee24862]
	Xét tính  tăng, giảm và bị chặn của dãy số $(u_n)$, biết $u_n=\dfrac{1}{\sqrt{1+n+n^2}}$
	 \choice
		 {Dãy số tăng, bị chặn trên}
		{Dãy số tăng,  bị chặn dưới}
		{\True  Dãy số giảm, bị chặn}
		{Dãy số không  bị chặn }
	 \loigiai{
		Ta  có $u_n>0, \,\forall n\in \mathbb{N}^*$ \\
		$\dfrac{u_{n+1}}{u_n}=\dfrac{\sqrt{n^2+n+1}}{\sqrt{(n+1)^2+(n+1)+1}}=\sqrt{\dfrac{n^2+n+1}{n^2+3n+3}}<1, \,\forall n \in  \mathbb{N}^*$.\\
		$ \Rightarrow  u_{n+1}<u_n, \,\forall n \in \mathbb{N}^* \Rightarrow $ dãy $(u_n)$ là dãy số giảm. \\
	Mặt khác $0<u_n<1 \Rightarrow $  dãy $(u_n)$ là dãy bị chặn.}
\end{ex}

% nguồn: Ok_TracNghiemLop11_Lan1
\begin{ex}%[1D3B2-4]%[ma:12641bb3]
	Xét tính bị  chặn của dãy số $(u_n)$ với $u_n=\dfrac{n^2+n+1}{n^2-n+1}$.
	 \choice
		 {\True Bị chặn}
		{Không  bị chặn}
		 {Bị chặn trên}
		{Bị chặn dưới} 
	\loigiai{ 
	Ta có $\forall n \in \mathbb{N}^*, 1<u_n<2 \Rightarrow  (u_n)$ bị chặn.}
\end{ex}

% nguồn: Ok_TracNghiemLop11_Lan1
\begin{ex}%[1D3B2-4]%[ma:212f4179]
	Xét tính bị chặn của dãy số $(u_n)$ với $u_n=\dfrac{1}{1.3}+\dfrac{1}{2.4}+\cdots +\dfrac{1}{n.(n+2)}$. 
	\choice 
		{\True Bị  chặn}
		{Không bị chặn} 
		{Bị  chặn trên}
		{Bị chặn  dưới}
	 \loigiai{
	Ta có $\forall n \in \mathbb{N}^*, 0<u_n<\dfrac{1}{1.2}+\dfrac{1}{2.3}+\cdots +\dfrac{1}{n.(n+1)}=1-\dfrac{1}{n+1}<1  \Rightarrow $ dãy $(u_n)$ bị chặn.}
\end{ex}

% nguồn: Ok_TracNghiemLop11_Lan1
\begin{ex}%[1D3B2-4]%[ma:24c5f850]
	Xét tính bị chặn của dãy số  $(u_n)$ với $u_n=3n-1$.
	\choice 
		{Bị  chặn}
		{Không  bị chặn}
		{Bị chặn trên} 
		{\True Bị  chặn dưới}
	 \loigiai{
	Ta có  $u_n \geq 2 \,\forall n \Rightarrow (u_n)$ bị chặn dưới; dãy $(u_n)$ không bị chặn trên.}
\end{ex}

% nguồn: Ok_TracNghiemLop11_Lan1
\begin{ex}%[1D3B2-4]%[ma:9b070ca3]
	 Xét tính bị chặn của các dãy số sau: $u_n=\dfrac{2n+1}{n+2}$.
	\choice 
		{\True Bị chặn} 
		{Không bị chặn} 
		{Bị chặn trên} 
		{Bị chặn  dưới}
	\loigiai{ 
	Ta có $\forall n \in \mathbb{N}^*, 0<u_n<2$ nên  dãy $(u_n)$ bị chặn.}
\end{ex}

% nguồn: Ok_TracNghiemLop11_Lan1
\begin{ex}%[1D3B2-4]%[ma:61fa1d5a]
	Xét  tính bị chặn của dãy số $(u_n)$ với $u_n=\dfrac{1}{1.3}+\dfrac{1}{3.5}+\cdots +\dfrac{1}{(2n-1)(2n+1)}$.
	\choice 
		{\True Bị chặn} 
		 {Không bị chặn}
		 {Bị chặn trên}
		{Bị chặn  dưới}
	 \loigiai{
	Ta có $\forall n \in \mathbb{N}^*,  u_n=\dfrac{n}{2n+1} \Rightarrow 0<u_n<1 \Rightarrow $ dãy $(u_n)$ bị chặn.}
\end{ex}

% nguồn: Ok_TracNghiemLop11_Lan1
\begin{ex}%[1D3K2-4]%[ma:9eb261c1]
	Cho dãy số $(x_n) \colon \left\{\begin{aligned}
	&x_0=1  \\
	& x_n=\dfrac{2n}{(n-1)^2}\cdot\displaystyle\sum_{i=1}^{n-1} x_i,  n=2,3,\cdots \\
\end{aligned}\right.$.\\
Khẳng định nào đúng về dãy  $(y_n)$ với $y_n=x_{n+1}-x_n$.
\choice 
	{\True Tăng,  bị chặn}
	{Giảm, bị  chặn}
	{Tăng, chặn  dưới}
	 {Giảm, chặn trên}
 \loigiai{
	Ta  có\\
	\begin{align*}
		x_{n+1}&=\dfrac{2(n+1)}{n^2}\displaystyle\sum_{i=1}^n x_i=\dfrac{2(n+1)}{n^2}\left(x_n+\displaystyle\sum_{i=1}^{n-1} x_i\right) \\ 
		 &=\dfrac{2(n+1)}{n^2}\left(x_n+\dfrac{(n-1)^2}{2n}x_n\right)=\dfrac{(n+1)(n^2+1)}{n^3}x_n.
	\end{align*}
	Do đó: $y_n=x_{n+1}-x_n=\dfrac{n^2+n+1}{n^3}x_n$ \\ 
	$\bullet $ Ta chứng minh dãy  $(y_n)$ tăng. \\
	 Ta có $y_{n+1}-y_n=\dfrac{(n+1)^2+n+2}{(n+1)^3}.\dfrac{(n+1)(n^2+1)}{n^3}x_n-\dfrac{n^2+n+1}{n^3}x_n$ \\
	$=\dfrac{(n^2+3n+3)(n^2+1)-(n^2+n+1)(n^2+2n+1)}{n^3(n+1)^2}x_n$  \\
	 $=\dfrac{2x_n}{n^3(n+1)^2}>0$, $\,\forall n=1,2$,.. \\
	$\bullet $ Ta chứng minh dãy $(y_n)$ bị  chặn. \\
	Trước hết ta chứng minh:  $x_n \leq 4(n-1)$ (1) với $\,\forall n=2,3\cdots $ \\
	* Với $n=2$, ta có $x_2=4x_1=4$ nên (1)  đúng với $n=2$ \\
	* Giả sử (1)  đúng với $n$, tức là: $x_n \leq 4(n-1)$, ta có \\
	 $x_{n+1}=\dfrac{(n+1)(n^2+1)}{n^3}x_n \leq \dfrac{4(n^4-1)}{n^3}<4n$ \\
	Nên (1) đúng với $n+1$. Theo nguyên lí quy  nạp ta suy ra (1) đúng \\
	Ta có $y_n=\dfrac{n^2+n+1}{n^3}x_n \leq  \dfrac{4(n-1)(n^2+n+1)}{n^3}=\dfrac{4(n^3-1)}{n^3}<4$ \\
Vậy bài toán được chứng minh.} 
\end{ex}

% nguồn: Ok_TracNghiemLop11_Lan1
\begin{ex}%[1D3K2-4]%[ma:118f23ec]
	Cho dãy số $(u_n)$, biết: $u_n=1+\dfrac{1}{2^2}+\dfrac{1}{3^2}+\cdots  +\dfrac{1}{n^2}$. Gọi $M, m$ sao cho $m\le u_n\le M,\,\forall n \in \mathbb{N}^*$. Chọn khẳng định \textbf{đúng}?
	 \choice
		{\True $m, M  \in (0;4)$}
		{$m, M \in  (1;4)$}
		 {$m, M \in (0;2)$}
		{$m, M \in (1;5)$} 
	\loigiai{ 
		Ta có  $u_{n+1}-u_n=\dfrac{1}{(n+1)^2}>0 \Rightarrow $ dãy $(u_n)$ là dãy số tăng. \\
		Do $u_n<1+\dfrac{1}{1.2}+\dfrac{1}{2.3}+\cdots  +\dfrac{1}{(n-1)n}=2+\dfrac{1}{n}$ \\
	$ \Rightarrow 1<u_n<3 \,\forall n  \geq 1 \Rightarrow $ dãy $(u_n)$ là dãy bị chặn.}
\end{ex}

% nguồn: Ok_TracNghiemLop11_Lan1
\begin{ex}%[1D3B2-4]%[ma:e25a31a7]
	Xét tính tăng giảm và bị chặn dãy số $(u_n)$ với $(u_n) \colon u_n=\dfrac{n+1}{n+2}$. 
	 \choice
		{\True Tăng,  bị chặn}
		 {Giảm, bị chặn}
		{Tăng, chặn dưới} 
		{Giảm, chặn trên} 
	\loigiai{ 
		Ta có $\forall  n \in \mathbb{N}^*,, u_{n+1}-u_n=\dfrac{n+2}{n+3}-\dfrac{n+1}{n+2}=\dfrac{(n+2)^2-(n+3)(n+1)}{(n+2)(n+3)}=\dfrac{1}{(n+2)(n+3)}>0$. \\
		 Mặt khác $\forall n \in \mathbb{N}^*, u_n=1-\dfrac{1}{n+2} \Rightarrow 0<u_n<1$. \\
	Vậy  dãy $(u_n)$ là dãy tăng và bị chặn.}
\end{ex}

% nguồn: Ok_TracNghiemLop11_Lan1
\begin{ex}%[1D3K2-4]%[ma:8f2d43b7]
	Cho dãy số $\left(u_n\right)$ cho bởi $u_n=\dfrac{1}{1.2}+\dfrac{1}{2.3}+\cdots +\dfrac{1}{n(n+1)},\,\forall n \in \mathbb{N}^*$.  Gọi $M, m$ sao cho \break $m\le u_n\le M,\,\forall n \in \mathbb{N}^*$. Tính $M+m$.
	\choice 
		{$\dfrac{2}{3}$} 
		 {$\sqrt{3}$}
		{\True  $\dfrac{3}{2}$}
		{$\dfrac{1}{2}$} 
	\loigiai{ 
		Ta có $\forall n \in \mathbb{N}^*, u_n=\dfrac{1}{1.2}+\dfrac{1}{2.3}+\cdots +\dfrac{1}{n(n+1)}\ge  \dfrac{1}{2}$.\\
		Mặt khác $\forall n \in  \mathbb{N}^*, u_n=1-\dfrac{1}{2}+\dfrac{1}{2}-\dfrac{1}{3}+\cdots +\dfrac{1}{n-1}-\dfrac{1}{n}+\dfrac{1}{n}-\dfrac{1}{n+1}=1-\dfrac{1}{n+1}\le 1 \Rightarrow M=1$.\\
	Vậy $m=\dfrac{1}{2},  M=1 \Rightarrow M+m=\dfrac{3}{2}$.}
\end{ex}

% nguồn: Ok_TracNghiemLop11_Lan1
\begin{ex}%[1D3B2-4]%[ma:5920a3a2]
	Xét tính bị chặn của dãy số $(u_n)$ với  $u_n=(-1)^n$.
	\choice 
		{\True Bị  chặn}
		 {Không bị chặn}
		{Bị chặn trên} 
		{Bị chặn dưới} 
	 \loigiai{
	Ta có  $\forall n \in \mathbb{N}^*, -1 \leq u_n \leq 1 \Rightarrow (u_n)$ là dãy bị chặn.}
\end{ex}

% nguồn: Ok_TracNghiemLop11_Lan1
\begin{ex}%[1D3B2-4]%[ma:9efac565]
	Cho dãy số $(u_n)$, biết $u_n=\dfrac{2^n}{n!}$. Khẳng định nào  sau đây \textbf{đúng}?
	\choice 
		{Dãy số tăng, bị  chặn trên}
		{Dãy  số tăng, bị chặn dưới}
		{ \True Dãy số giảm, bị chặn trên} 
		{Dãy số giảm, bị chặn dưới} 
	 \loigiai{
		 Ta có $\dfrac{u_{n+1}}{u_n}=\dfrac{2^{n+1}}{(n+1)!} \colon \dfrac{2^n}{n!}=\dfrac{2^{n+1}}{(n+1)!}.\dfrac{n!}{2^n}=\dfrac{2}{n+1}<1, \forall n \in \mathbb{N}^*$ \\
		 Mà $\forall n \in \mathbb{N}^*, u_n>0, \Rightarrow u_{n+1}<u_n \Rightarrow $ dãy $(u_n)$ là dãy số giảm. \\
	Vì $0<u_n \leq u_1=2, \,\forall n \geq 1 \Rightarrow $ dãy $(u_n)$  là dãy bị chặn.}
\end{ex}


%%%%%%%%%%%%%%%% ĐÚNG/SAI %%%%%%%%%%%%%%%%

% nguồn: Ok_dungsai_lop11_Lan1
\begin{ex}%[1D3N2-4]%[ma:98fa0747]
  Mỗi kết quả giới hạn sau đúng hay sai?
  \choiceTF
  {\True $\lim\limits_{x \to  0} \dfrac{1}{\left|x\right|}=+\infty$}
  {$\lim\limits_{x \to  1} \dfrac{x}{\left|x^2-x\right|}=\dfrac{1}{2}$}
  {\True $\lim\limits_{x  \to 2} \dfrac{x-1}{\left|x-2\right|}=-\infty$}
  {\True $\lim\limits_{x \to 0}  \dfrac{\left|x^2-x\right|}{x^2}=+\infty$}
 \loigiai{ 
  \begin{itemchoice}
   \itemch $\lim\limits_{x \to 0} \dfrac{1}{\left|x\right|}=+\infty$, vì $\lim\limits_{x \to 0} \left|x\right|=0$ và $\left|x\right|>0$ khi $x\ne 0$. 
   \itemch $\lim\limits_{x \to 1} \dfrac{x}{\left|x^2-x\right|}=\lim\limits_{x \to 1} \dfrac{1}{\left|x-1\right|}=+\infty$, vì $\lim\limits_{x \to 1} \left|x-1\right|=0$ và $\left|x-1\right|>0$  khi $x\ne 1$.
   \itemch $\lim\limits_{x \to 2}  \dfrac{1-x}{\left|x-2\right|}=-\infty
   $, vì $\lim\limits_{x \to  2} (1-x)=-1<0$, $\lim\limits_{x \to 2} \left|x-2\right|=0$ và $\left|x-2\right|>0$ khi $x\ne 2$.
   \itemch $\lim\limits_{x \to 0} \dfrac{\left|x^2-x\right|}{x^2}=\lim\limits_{x \to 0} \dfrac{\left|x-1\right|}{\left|x\right|}=+\infty$, vì $\lim\limits_{x \to 0} \left|x-1\right|=1>0$, $\lim\limits_{x \to 0} \left|x\right|=0$ và $\left|x\right|>0$ khi  $x \ne 0$.
  \end{itemchoice}
  }
\end{ex}

% nguồn: Ok_dungsai_lop11_Lan1
\begin{ex}%[1D3H2-4]%[ma:cb118ed3]
 Mỗi kết quả giới hạn sau đúng hay sai? 
 \choiceTF 
  {\True $\lim\limits_{x \to 1}  \dfrac{\left|x^2-x\right|}{(x-1)^2}=+\infty$}
  {$\lim\limits_{x \to 0}  \dfrac{x}{x^2}=1$}
  {$\lim\limits_{x \to -2} \dfrac{x^2+4x+4}{x^2+3x+2}=-\infty$} 
   {\True $\lim\limits_{x \to 3} \left(\dfrac{1}{x}-\dfrac{1}{3}\right)\dfrac{1}{(x-3)^3}=-\infty$}
  \loigiai{
  \begin{itemchoice}
    \itemch $\lim\limits_{x \to 1} \dfrac{\left|x^2-x\right|}{(x-1)^2}
   =\lim\limits_{x \to 1}  \dfrac{\left|x\right|}{\left|x-1\right|}
   =+\infty$, vì $\lim\limits_{x \to 1} \left|x\right|=1$, $\lim\limits_{x \to  1} \left|x-1\right|=0$ và $\left|x-1\right|>0$ với mọi $x\ne 1$.
   \itemch $\lim\limits_{x \to  0} \dfrac{x}{x^2}=\lim\limits_{x \to 0} \dfrac{1}{x}$.
   \begin{itemize}
    \item  Với $x>0$, ta được $\lim\limits_{x \to 0} \dfrac{1}{x}=+\infty$;
    \item Với  $x<0$, ta được $\lim\limits_{x \to 0} \dfrac{1}{x}=-\infty$.
   \end{itemize}
   Do đó không tồn tại giới hạn  $\lim\limits_{x \to 0} \dfrac{x}{x^2}$.
   \itemch $\lim\limits_{x \to  -2} \dfrac{x^2+4x+4}{x^2+3x+2}
   =\lim\limits_{x \to -2} \dfrac{(x+2)^2}{(x+2)(x+1)} 
   =\lim\limits_{x  \to -2} \dfrac{x+2}{x+1}=0$.
   \itemch $\lim\limits_{x  \to 3} \left(\dfrac{1}{x}-\dfrac{1}{3}\right)\dfrac{1}{(x-3)^3}
   =\lim\limits_{x \to 3} \left(\dfrac{3-x}{3x}\cdot  \dfrac{1}{(x-3)^3}\right)
   =\lim\limits_{x  \to 3} \left(-\dfrac{1}{3x}\cdot \dfrac{1}{(x-3)^2}\right)=-\infty$;\\
   Vì $\lim\limits_{x \to 3} \left(-\dfrac{1}{3x}\right)=-\dfrac{1}{9}<0$  và $\lim\limits_{x \to 3} \dfrac{1}{(x-3)^2}=+\infty$
  \end{itemchoice}
  }
\end{ex}

% nguồn: Ok_dungsai_lop11_Lan1
\begin{ex}%[1D3H2-4]%[ma:128c86cc]
 Tìm  được các giới hạn sau.
 \choiceTF 
  {\True $\lim\limits_{x\to{2^+}}(\sqrt{x+2}-1)=1$} 
  {\True $\lim\limits_{x\to{1^+}}\dfrac{4x-3}{x-1}=+\infty  $}
  {\True $\lim\limits_{x\to{2^-}}\left(\dfrac{1}{x-2}-\dfrac{1}{x^2-4}\right)=-\infty $} 
   {$\lim\limits_{x\to-1^-}\dfrac{|x+1|}{x^2-1}=-\infty$}
  \loigiai{
  \begin{itemchoice}
    \itemch $\lim\limits_{x\to{2^+}}(\sqrt{x+2}-1)=\sqrt{2+2}-1=1$.
   \itemch $\lim\limits_{x\to{1^+}}\dfrac{4x-3}{x-1}= \lim\limits_{x\to{1^+}}\left[(4x-3)\cdot\dfrac{1}{x-1}\right]=+\infty$  \\
   vì $\lim\limits_{x\to{1^+}}(4x-3)=1$, $\lim\limits_{x\to{1^+}}\dfrac{1}{x-1}=+\infty $. 
   \itemch \begin{eqnarray*}
    \lim\limits_{x\to{2^-}}\left(\dfrac{1}{x-2}-\dfrac{1}{x^2-4}\right)&=& \lim\limits_{x\to{2^-}}\dfrac{x+2-1}{(x-2)(x+2)}\\
    &=& \lim\limits_{x\to{2^-}}\dfrac{x+1}{(x-2)(x+2)}\\ &=&\lim\limits_{x\to{2^-}}\left[\dfrac{x+1}{x+2}\cdot\dfrac{1}{(x-2)}\right]= -\infty,
  \end{eqnarray*}
   do $\heva{
  &\lim\limits_{x\to{2^-}}\dfrac{x+1}{x+2}=\dfrac{3}{4}\\ 
  &\lim\limits_{x\to{2^-}}\dfrac{1}{x-2}=-\infty}$ 
  \itemch  $\lim\limits_{x\to-1^-}\dfrac{|x+1|}{x^2-1}=\lim\limits_{x\to-1^-}\dfrac{-x-1}{(x-1)(x+1)}=\lim\limits_{x\to-1^-}\dfrac{-1}{x-1}=\dfrac{1}{2}$.
 \end{itemchoice}
} 
\end{ex}

% nguồn: Ok_dungsai_lop11_Lan2
\begin{ex}%[1D3H2-4]%[ma:9e5b5928]
	Cho $f(x)=\sqrt{4 x^2-7 x+12}$ và  $g(x)=a|x|-17$ với $a$ là tham số. Đặt $L=\lim\limits_{x \to -\infty}\dfrac{f(x)}{g(x)}$.
	 \choiceTF
		{\True Với $a=-1$ thì $\lim\limits_{x \to  -\infty}g(x)=-\infty$}
		{\True $\lim\limits_{x \to  -\infty}f(x)=+\infty$}
		{$\lim\limits_{x  \to -\infty}\dfrac{f(x)}{x}=2$}
		 {\True Nếu $L=\dfrac{2}{3}$ thì $a=3$}
	 \loigiai{
		\begin{itemchoice}
			\itemch \textbf{Đúng.} Với $a=-1$ thì  $\lim\limits_{x \to -\infty}g(x)=\lim\limits_{x \to -\infty}(x-17)=-\infty$.
			\itemch \textbf{Đúng.} Ta  có
			\begin{align*}
				 \lim\limits_{x \to -\infty}f(x)
				&=\lim\limits_{x \to -\infty}\sqrt{4 x^2-7 x+12} \\ 
				&=\lim\limits_{x  \to -\infty}|x|\sqrt{4-\dfrac{7}{x}+\dfrac{12}{x^2}} \\
				&=-\lim\limits_{x  \to -\infty}x \sqrt{4-\dfrac{7}{x}+\dfrac{12}{x^2}}=+\infty.
			\end{align*}
			\itemch \textbf{Sai.} Ta thấy 
			 \[
			\lim\limits_{x  \to -\infty}\dfrac{f(x)}{x}
			=\lim\limits_{x  \to -\infty}\dfrac{-x \sqrt{4-\dfrac{7}{x}+\dfrac{12}{x^2}}}{x}
			=-\lim\limits_{x  \to -\infty}\sqrt{4-\dfrac{7}{x}+\dfrac{12}{x^2}}
			 =-\sqrt{{4}}=-2.
			\] 
			\itemch \textbf{Đúng.} Ta có 
			\begin{align*}
				 \lim\limits_{x \to -\infty}\dfrac{\sqrt{4 x^2-7 x+12}}{a|x|-17}
				&=\lim\limits_{x  \to -\infty}\dfrac{-x \sqrt{4-\dfrac{7}{x}+\dfrac{12}{x^2}}}{-x\left(a+\dfrac{17}{x}\right)} \\
				&=\lim\limits_{x \to  -\infty}\dfrac{\sqrt{4-\dfrac{7}{x}+\dfrac{12}{x^2}}}{a+\dfrac{17}{x}} \\
				 &=\dfrac{2}{a}=\dfrac{2}{3}.
			\end{align*}
			 Suy ra $a=3$.
		\end{itemchoice}
	} 
\end{ex}

% nguồn: Ok_dungsai_lop11_Lan1
\begin{ex}%[1D3V2-4]%[ma:23cf3444]
 Tìm được các giới hạn  sau.
 \choiceTF 
  {\True $\lim\limits_{x\to-\infty}\left(x^2-10x\right)=+\infty $} 
  {\True $\lim\limits_{x\to+\infty}\dfrac{3x^2-4x+1}{2x^2+x+1}=\dfrac{3}{2}$} 
   {$\lim\limits_{x\to-\infty}\dfrac{\sqrt{x^2+x+1}-3x}{2-3x}=\dfrac{5}{4}$}
   {\True $\lim\limits_{x\to-\infty}\dfrac{\sqrt[3]{8x^3+3x^2+1}-x}{\sqrt{4x^2-x+2}+3x}=1$}
 \loigiai{ 
  \begin{itemchoice}
   \itemch $\lim\limits_{x\to-\infty}\left(x^2-10x\right)= \lim\limits_{x\to-\infty}{x^2}\left(1-\dfrac{10}{x}\right)=+\infty $. 
   \itemch  $\lim\limits_{x\to+\infty}\dfrac{3x^2-4x+1}{2x^2+x+1}= \lim\limits_{x\to+\infty}\dfrac{x^2\left(3-\dfrac{4}{x}+\dfrac{1}{x^2}\right)}{x^2\left(2+\dfrac{1}{x}+\dfrac{1}{x^2}\right)}= \lim\limits_{x\to+\infty}\dfrac{3-\dfrac{4}{x}+\dfrac{1}{x^2}}{2+\dfrac{1}{x}+\dfrac{1}{x^2}}=\dfrac{3}{2}$
   \itemch \begin{eqnarray*}
    \lim\limits_{x\to-\infty}\dfrac{\sqrt{x^2+x+1}-3x}{2-3x}&=& \lim\limits_{x\to-\infty}\dfrac{\sqrt{x^2\left(1+\dfrac{1}{x}+\dfrac{1}{x^2}\right)}-3x}{x\left(\dfrac{2}{x}-3\right)}\\
   &=& \lim\limits_{x\to-\infty}\dfrac{-x\sqrt{1+\dfrac{1}{x}+\dfrac{1}{x^2}}-3x}{x\left(\dfrac{2}{x}-3\right)}\\ 
   &=&\lim\limits_{x\to-\infty}\dfrac{-\sqrt{1+\dfrac{1}{x}+\dfrac{1}{x^2}}-3}{\dfrac{2}{x}-3}=\dfrac{-\sqrt  1-3}{-3}=\dfrac{4}{3}.
  \end{eqnarray*}
  \itemch \begin{eqnarray*}
   \lim\limits_{x\to-\infty}\dfrac{\sqrt[3]{8x^3+3x^2+1}-x}{\sqrt{4x^2-x+2}+3x}&=&\lim\limits_{x\to-\infty}\dfrac{\sqrt[3]{x^3\left(8+\dfrac{3}{x}+\dfrac{1}{x^3}\right)}-x}{\sqrt{x^2\left(4-\dfrac{1}{x}+\dfrac{2}{x^2}\right)}+3x}\\
  &=&\lim\limits_{x\to-\infty}\dfrac{x\sqrt[3]{8+\dfrac{3}{x}+\dfrac{1}{x^3}}-x}{x\sqrt{4-\dfrac{1}{x}+\dfrac{2}{x^2}}+3x}\\ 
  &=&\lim\limits_{x\to-\infty}\dfrac{\sqrt[3]{8+\dfrac{3}{x}+\dfrac{1}{x^3}}-1}{\sqrt{4-\dfrac{1}{x}+\dfrac{2}{x^2}}+3}= \dfrac{\sqrt[3]{8}-1}{-\sqrt  4+3}=1.
 \end{eqnarray*}
 \end{itemchoice}
 }
\end{ex}


%%%%%%%%%%%%%%%% TỰ LUẬN %%%%%%%%%%%%%%%%

% nguồn: Ok_TuanLuan_Toan11_Lan2
\begin{ex}%[1D3H2-4]%[ma:58ac804d]
	Giới hạn $I=\lim\limits_{x\to+\infty}\left(-2x^3+4x+5\right)$  bằng
	\loigiai{ 
		Ta có  $I=\lim\limits_{x\to+\infty} \left(-2x^3+4x+5\right)=\lim\limits_{x\to+\infty} \left[x^3\left(-2+\dfrac{4}{x^2}+\dfrac{5}{x^3} \right)\right]=-\infty$.\\
		Vì $\lim\limits_{x\to+\infty}x^3=+\infty$,  $\lim\limits_{x\to+\infty} \left(-2+\dfrac{4}{x^2}+\dfrac{5}{x^3}\right)=-2<0$.
	 }
\end{ex}

% nguồn: Ok_TuanLuan_Toan11_Lan2
\begin{ex}%[1D3H2-4]%[ma:53fd54e4]
	 Tìm tất cả các giá trị của $a$ để $\lim\limits_{x\to -\infty}\left(\sqrt{2x^2+1}+ax\right)$ là $+\infty .$
	\loigiai{ 
		Vì $\lim\limits_{x\to -\infty}x=-\infty$  nên $\lim\limits_{x\to -\infty}\left(\sqrt{2x^2+1}+ax\right)=\lim\limits_{x\to -\infty}x\left(-\sqrt{2+\dfrac{1}{x^2}}+a\right)=+\infty$ khi và chỉ khi
		 $$\lim\limits_{x\to -\infty}\left(-\sqrt{2+\dfrac{1}{x^2}}+a\right)=a-\sqrt 2<0\Leftrightarrow a<\sqrt 2.$$
	} 
\end{ex}

% nguồn: Ok_TuanLuan_Toan11_Lan2
\begin{ex}%[1D3H2-4]%[ma:613646a2]
	Tính  giới hạn $\lim\limits_{x\to-\infty}\dfrac{x^2+1}{x-2}$.
	\loigiai{ 
		$\lim\limits_{x\to-\infty} \dfrac{x^2+1}{x-2}=\lim\limits_{x\to-\infty} \dfrac{x^2}{x}\cdot  \dfrac{1+\dfrac{1}{x^2}}{1-\dfrac{2}{x}}=\lim\limits_{x\to-\infty} x\cdot \dfrac{1+\dfrac{1}{x^2}}{1-\dfrac{2}{x}}$.\\
		Do  $\lim\limits_{x\to-\infty}x=-\infty$ và $\lim\limits_{x\to-\infty} \dfrac{1+\dfrac{1}{x^2}}{1-\dfrac{2}{x}}=1$ nên $\lim\limits_{x\to-\infty}\dfrac{x^2+1}{x-2}=-\infty$.
	 }
\end{ex}

% nguồn: Ok_TuanLuan_Toan11_Lan2
\begin{ex}%[1D3H2-4]%[ma:8b991f69]
	Cho hàm số $f(x)=\heva{&\dfrac{x^2+1}{1-x}&&\text{với}~ x<1\\&\sqrt{2x-2}&&\text{với}~ x\ge 1}$. Khi đó $\lim\limits_{x\to 1^{-}} f(x)$ là 
	 \loigiai{
		$\lim\limits_{x\to 1^{-}} f(x)=\lim\limits_{x\to 1^{-}} \dfrac{x^2+1}{1-x}=+\infty$  vì $\heva{&\lim\limits_{x\to 1^{-}} \left(x^2+1\right)=2\\&\lim\limits_{x\to 1^{-}}(1-x)=0\\& 1-x> 0,~\forall x<1.}$
	} 
\end{ex}

% nguồn: Ok_TuanLuan_Toan11_Lan2
\begin{ex}%[1D3H2-4]%[ma:649ebfa5]
	$\lim\limits_{x\to+\infty}\left(-2x^3-2x\right)$ bằng 
	\loigiai{ 
		Ta có $\lim\limits_{x\to+\infty} (-2x^3-2x)=\lim\limits_{x\to+\infty} x^3\left(-2-\dfrac{2}{x^2}\right)$.\\ 
		Mà  $\lim\limits_{x\to+\infty}x^3=+\infty$; $\lim\limits_{x\to+\infty} \left(-2-\dfrac{2}{x^2}\right)=-2<0$ nên $\lim\limits_{x\to+\infty} x^3\left(-2-\dfrac{2}{x^2}\right)=-\infty$.\\
		 Vậy $\lim\limits_{x\to+\infty}\left(-2x^3-2x\right)=-\infty$.
	 }
\end{ex}

% nguồn: Ok_TuanLuan_Toan11_Lan2
\begin{ex}%[1D3H2-4]%[ma:b9e21447]
	Có bao nhiêu giá trị nguyên của tham  số $m$ thuộc đoạn $[-5;5]$ để\\ $\lim\limits_{x\to+\infty}\left[x-2\left(m^2-4\right)x^3\right]=-\infty$.
	 \loigiai{
		Ta  có $\lim\limits_{x\to+\infty} \left[x-2\left(m^2-4\right)x^3\right]=\lim\limits_{x\to+\infty} x^3\left[\dfrac{1}{x^2}-2\left(m^2-4\right)\right]$.\\
		Ta  có $\lim\limits_{x\to+\infty}x^3=+\infty$ và $\lim\limits_{x\to+\infty}\left[\dfrac{1}{x^2}-2\left(m^2-4\right)\right]=-2\left(m^2-4\right)$ nên\\ $\lim\limits_{x\to+\infty} \left[x-2\left(m^2-4\right)x^3\right]=-\infty \Leftrightarrow-2(m^2-4)<0\Leftrightarrow m^2-4> 0\Leftrightarrow \hoac{&m>2\\&m <-2.}$\\
		 Lại có $m$ thuộc đoạn $[-5;5]$ nên các giá trị nguyên thỏa mãn bài toán của $m$ là $\{-5;-4;-3;3;4;5\}$.\\
		Vậy có  $6$ số nguyên thỏa mãn bài toán.
	 }
\end{ex}

% nguồn: Ok_TuanLuan_Toan11_Lan2
\begin{ex}%[1D3H2-4]%[ma:72eecc60]
	Cho hàm số $f(x)=\heva{&x^2-3 &&\text{với}~ x\ge 2\\&x-1&& \text{với}~x<2}$. Khi đó  $\lim\limits_{x\to 2} f(x)$ là
	 \loigiai{
		Ta có $\heva{&\lim\limits_{x\to 2^{+}}f(x)=\lim\limits_{x\to 2^{+}} \left(x^2-3\right)=1\\&\lim\limits_{x\to 2^{-}} f(x)=\lim\limits_{x\to 2^{-}}(x-1)=1} \Rightarrow \lim\limits_{x\to 2^{+}} f(x)=\lim\limits_{x\to  2^{-}} f(x)=1\Rightarrow \lim\limits_{x\to 2} f(x)=1$.
	} 
\end{ex}

% nguồn: Ok_TuanLuan_Toan11_Lan2
\begin{ex}%[1D3H2-4]%[ma:294687d6]
	Cho hàm số $f(x)=\heva{&\dfrac{2x}{\sqrt{1-x}}&&\text{với}~ x<1\\ &\sqrt{3x^2+1}&& \text{với}~ x\ge 1}$. Khi đó $\lim\limits_{x\to 1^{+}}f(x)$  là
	\loigiai{ 
		$\lim\limits_{x\to 1^{+}} f(x)=\lim\limits_{x\to  1^{+}} \sqrt{3x^2+1}=\sqrt{3\cdot 1^2+1}=2$.
	} 
\end{ex}

% nguồn: Ok_TuanLuan_Toan11_Lan2
\begin{ex}%[1D3H2-4]%[ma:225a172c]
	Tính giới  hạn $\lim\limits_{x\to-\infty}\left(2x^3-x^2+1\right)$.
	 \loigiai{
		Ta có $\lim\limits_{x\to-\infty}  \left(2x^3-x^2+1\right)=\lim\limits_{x\to-\infty} \left[x^3\left(2-\dfrac{1}{x}+\dfrac{1}{x^3}\right)\right]=-\infty$.\\ Vì $\lim\limits_{x\to-\infty} x^3=-\infty$ và $\lim\limits_{x\to-\infty} \left(2-\dfrac{1}{x}+\dfrac{1}{x^3}\right)=2$.
	} 
\end{ex}

% nguồn: Ok_TuanLuan_Toan11_Lan2
\begin{ex}%[1D3H2-4]%[ma:dee63d9c]
	$\lim\limits_{x\to-\infty}\left(-3x^3+2x\right)$. 
	 \loigiai{
		Ta có $\lim\limits_{x\to-\infty} \left(-3x^3+2x\right)=\lim\limits_{x\to-\infty}  \left[x^3\left(-3+\dfrac{2}{x^2}\right)\right]$.\\
		 Ta thấy $\heva{&\lim\limits_{x\to-\infty}x^3=-\infty\\ &\lim\limits_{x\to-\infty}\left(-3+\dfrac{2}{x^2}\right)=-3< 0} \Rightarrow \lim\limits_{x\to-\infty}\left[x^3\left(-3+\dfrac{2}{x^2} \right)\right]=+\infty$.\\
		 Vậy $\lim\limits_{x\to-\infty}\left(-3x^3+2x\right)=+\infty$.
	 }
\end{ex}

% nguồn: Ok_TuanLuan_Toan11_Lan2
\begin{ex}%[1D3H2-4]%[ma:fef1f760]
	Kết quả của giới hạn  $\lim\limits_{x\to 2^{+}}\dfrac{x-15}{x-2}$ là
	\loigiai{ 
		 Vì $\heva{&\lim\limits_{x\to 2^{+}}(x-15)=-13<0\\&\lim\limits_{x\to 2^{+}}(x-2)=0\\& x-2>0,~\forall x>2}$ nên $\lim\limits_{x\to 2^{+}}\dfrac{x-15}{x-2}=-\infty$.
	} 
\end{ex}

% nguồn: Ok_TuanLuan_Toan11_Lan2
\begin{ex}%[1D3H2-4]%[ma:f204ba35]
	 Biết $\lim\limits_{x\to 1}f(x)=-2$. Khi đó $\lim\limits_{x\to 1} \dfrac{f(x)}{\left(x-1\right)^2}$ bằng
	 \loigiai{
		Có $\lim\limits_{x\to 1}f(x)=-2<0$, $\lim\limits_{x\to 1}  (x-1)^2=0$ và $(x-1)^2>0$, $\forall x\ne 1$ nên $\lim\limits_{x\to 1} \dfrac{f(x)}{(x-1)^2}=-\infty$.
	 }
\end{ex}

% nguồn: Ok_TuanLuan_Toan11_Lan2
\begin{ex}%[1D3H2-4]%[ma:15e58acd]
	Biết $\lim\limits_{x\to-2}f(x)=-1$. Khi  đó $\lim\limits_{x\to-2} \dfrac{f(x)}{(x+2)^4}$ bằng bao nhiêu?
	\loigiai{ 
		Ta  có $\lim\limits_{x\to-2}f(x)=-1<0$,
		$\lim\limits_{x\to-2}(x+2)^4=0$  và với $\forall x\ne-2$ thì $(x+2)^4>0$.\\
		Suy ra  $\lim\limits_{x\to-2}\dfrac{f(x)}{(x+2)^4}=-\infty$.
	 }
\end{ex}

% nguồn: Ok_TuanLuan_Toan11_Lan2
\begin{ex}%[1D3H2-4]%[ma:ced87a57]
	Kết quả của giới hạn $\lim\limits_{x\to  2^{-}} \dfrac{\left|2-x\right|}{2x^2-5x+2}$ là
	 \loigiai{
		 Ta có $\left|2-x\right|=2-x$, $\forall x<2$.\\
		$\lim\limits_{x\to  2^{-}} \dfrac{\left|2-x\right|}{2x^2-5x+2}=\lim\limits_{x\to 2^{-}} \dfrac{2-x}{\left(2-x\right)\left(1-2x\right)}=\lim\limits_{x\to 2^{-}} \dfrac{1}{1-2x}=-\dfrac{1}{3}$.
	 }
\end{ex}

% nguồn: Ok_TuanLuan_Toan11_Lan2
\begin{ex}%[1D3H2-4]%[ma:56da6456]
	 Kết quả của giới hạn $\lim\limits_{x\to 2^{+}} \dfrac{\sqrt{x+2}}{\sqrt{x-2}}$ là
	\loigiai{ 
		$\heva{&\lim\limits_{x\to 2^{+}}\sqrt{x+2}=2>0\\&\lim\limits_{x\to 2^{+}}\sqrt{x-2}=0\\&\sqrt{x-2}>0,~\forall x>2} \Rightarrow  \lim\limits_{x\to 2^{+}}\dfrac{\sqrt{x+2}}{\sqrt{x-2}}=+\infty$.
	} 
\end{ex}

% nguồn: Ok_TuanLuan_Toan11_Lan2
\begin{ex}%[1D3V2-4]%[ma:f30a21ad]
	Có bao nhiêu giá trị $m$ nguyên thuộc đoạn  $[-20;20]$ để\\ $\lim\limits_{x\to+\infty}\left(\sqrt{4x^2-3x+2}+mx-1\right)=-\infty$?
	\loigiai{ 
		 Ta có
		 \allowdisplaybreaks
		\begin{eqnarray*}
			 \lim\limits_{x\to+\infty} \left(\sqrt{4x^2-3x+2}+mx-1\right)&=&\lim\limits_{x\to+\infty} \left(x\sqrt{4-\dfrac{3}{x}+\dfrac{2}{x^2}}+mx-1\right)\\
			 &=&\lim\limits_{x\to+\infty} x\left(\sqrt{4-\dfrac{3}{x}+\dfrac{2}{x^2}}+m-\dfrac{1}{x}\right).
		\end{eqnarray*}
		Mà $\lim\limits_{x\to+\infty}x=+\infty$ và $\lim\limits_{x\to+\infty} \left(\sqrt{4-\dfrac{3}{x}+\dfrac{2}{x^2}}+m-\dfrac{1}{x}\right)=2+m$ nên  $\lim\limits_{x\to+\infty} \left(\sqrt{4x^2-3x+2}+mx-1\right)=-\infty$ khi $2+m<0\Leftrightarrow m <-2$.\\
		 Do $m$ nguyên thuộc đoạn $[-20;20]$ nên $m\in \{-20;-19;-18;\ldots;-3\}$.\\
		Vậy có $18$ giá trị $m$ nguyên thuộc đoạn $[-20;20]$ thỏa bài  toán.
	 }
\end{ex}

% nguồn: Ok_TuanLuan_Toan11_Lan2
\begin{ex}%[1D3V2-4]%[ma:0a265838]
	 Cho hàm số $f(x)=\heva{&\dfrac{\sqrt{x+4}-2}{x}&&\text{khi}~ x>0\\ &mx+m+\dfrac{1}{4}&&\text{khi}~ x\le 0}$, $m$ là tham số. Tìm giá trị của $m$ để hàm số có giới hạn tại $x=0$.
	\loigiai{ 
		Ta có 
		 \allowdisplaybreaks
		\begin{eqnarray*}
			\lim\limits_{x\to  0^{+}}f(x)&=&\lim\limits_{x\to 0^{+}} \dfrac{\sqrt{x+4}-2}{x}\\
			&=&\lim\limits_{x\to 0^{+}} \dfrac{\left(x+4\right)-2^2}{x\left(\sqrt{x+4}+2\right)}\\ 
			 &=&\lim\limits_{x\to 0^{+}} \dfrac{x}{x\left(\sqrt{x+4}+2\right)}\\
			 &=&\lim\limits_{x\to 0^{+}}\dfrac{1}{\sqrt{x+4}+2}=\dfrac{1}{4}.
		\end{eqnarray*}
		 $\lim\limits_{x\to 0^{-}}f(x)=\lim\limits_{x\to 0^{-}} \left(mx+m+\dfrac{1}{4}\right)=m+\dfrac{1}{4}$.\\
		Hàm số đã cho có giới hạn tại $x=0$ khi và chỉ  khi $\lim\limits_{x\to 0^{+}}f(x)=\lim\limits_{x\to 0^{-}}f(x)$.\\
		Do đó $m+\dfrac{1}{4}=\dfrac{1}{4}\Leftrightarrow  m=0$.
	} 
\end{ex}

% nguồn: Ok_TuanLuan_Toan11_Lan2
\begin{ex}%[1D3V2-4]%[ma:06ee05d0]
	Cho hàm số $f(x)=\heva{&\sqrt{x-2}+3&& \text{với}~ x\ge 2\\&ax-1 &&\text{với}~ x < 2}$. Tìm $a$ để tồn  tại $\lim\limits_{x\to 2} f(x)$.
	 \loigiai{
		Ta có $\heva{&\lim\limits_{x\to 2^{-}}f(x)=\lim\limits_{x\to 2^{-}} (ax-1)=2a-1\\&\lim\limits_{x\to  2^{+}}f(x)=\lim\limits_{x\to 2^{+}} \left(\sqrt{x-2}+3\right)=3.}$\\
		Khi đó $\lim\limits_{x\to 2}f(x)$ tồn tại  $\Leftrightarrow \lim\limits_{x\to 2^{-}}f(x)=\lim\limits_{x\to 2^{+}} f(x)\Leftrightarrow 2a-1=3\Leftrightarrow a=2$.
	 }
\end{ex}

% nguồn: Ok_TuanLuan_Toan11_Lan2
\begin{ex}%[1D3V2-4]%[ma:aa4135bc]
	Biết rằng $\dfrac{(2-a)x-3}{\sqrt{x^2+1}-x}$ có giới hạn là $+\infty $ khi $x\to +\infty $. Tính giá trị nhỏ  nhất của $P=a^2-2a+4$.
	\loigiai{ 
	 Ta có
	\begin{align*}
		\lim\limits_{x\to +\infty }\dfrac{(2-a)x-3}{\sqrt{x^2+1}-x}&=\lim\limits_{x\to +\infty  }\left[(2-a)x-3\right]\left( \sqrt{x^2+1}+x \right)\\
		 &=\lim\limits_{x\to +\infty }x^2\left(2-a-\dfrac{3}{x} \right)\left(\sqrt{1+\dfrac{1}{x^2}}+1 \right).
	\end{align*}
	Vì $\heva{& \lim\limits_{x\to +\infty }x^2=+\infty \\& \lim\limits_{x\to +\infty } \left( \sqrt{1+\dfrac{1}{x^2}}+1 \right)=4>0  }\Rightarrow \lim\limits_{x\to +\infty }\,\dfrac{(2-a)x-3}{\sqrt{x^2+1}-x}=+\infty $.\\
	Do đó $ \lim\limits_{x\to +\infty }\left(2-a-\dfrac{3}{x} \right)=2-a>0\Rightarrow a<2$. 
	 }
\end{ex}

% nguồn: Ok_TuanLuan_Toan11_Lan1
\begin{bt}%[1D3Y2-4]%[ma:c68d96dc]
	Cho dãy số  $(u_n)$, biết: $1 \le u_n\le 10, \forall n\in \mathbb{N}^*$. Xét tính bị chặn của dãy số.
	\loigiai{ 
	Vì $1 \le u_n\le 10, \forall n\in \mathbb{N}^*$ nên giá trị  lớn nhất của dãy số là $1$, giá trị lớn nhất của dãy số là $10$.}
\end{bt}

% nguồn: Ok_TuanLuan_Toan11_Lan1
\begin{bt}%[1D3B2-4]%[ma:7ee7a10f]
	Xét tính tăng, giảm và bị chặn của dãy số $(u_n)$, biết: $u_n=\dfrac{n^2+3n+1}{n+1}$. 
	\loigiai{ 
	Ta có $\forall n \in \mathbb{N}^*, u_n>\dfrac{n^2+2n+1}{n+1}=n+1 \geq 2  \Rightarrow $ dãy $(u_n)$ bị chặn dưới.}
\end{bt}

% nguồn: Ok_TuanLuan_Toan11_Lan2
\begin{ex}%[1D3Y2-4]%[ma:2689acdf]
	Cho dãy số $(u_n)$, biết: $1 \le u_n\le 10, \forall n\in \mathbb{N}^*$. Xét tính bị chặn của dãy số. 
	 \loigiai{
	Vì $1 \le u_n\le 10, \forall n\in \mathbb{N}^*$ nên giá trị lớn nhất của dãy số là $1$, giá trị lớn nhất của dãy số là  $10$.}
\end{ex}

% nguồn: Ok_TuanLuan_Toan11_Lan2
\begin{ex}%[1D3B2-4]%[ma:70d8c7c7]
	Xét tính  tăng, giảm và bị chặn của dãy số $(u_n)$, biết: $u_n=\dfrac{n^2+3n+1}{n+1}$.
	 \loigiai{
	Ta có $\forall n \in \mathbb{N}^*, u_n>\dfrac{n^2+2n+1}{n+1}=n+1 \geq 2 \Rightarrow  $ dãy $(u_n)$ bị chặn dưới.}
\end{ex}
